% !TeX TS-program = xelatex

\documentclass[compress]{beamer}

\beamertemplatenavigationsymbolsempty
\setbeamertemplate{headline}{}

% ==== Packages section ============================

\usepackage{xltxtra}
\usepackage[main=russian,english]{babel}

\setmainfont{DejaVu Serif}
\setsansfont{DejaVu Sans}
\setmonofont{Fira Mono}

\usepackage[export]{adjustbox}
\usepackage{tikz}

\usepackage[newfloat]{minted}
\usemintedstyle{borland}
\SetupFloatingEnvironment{listing}{name=Листинг}

\usepackage{fancyvrb}

\usepackage{hyperref}
\hypersetup{
    colorlinks=true,
    linkcolor=blue,
    filecolor=magenta,
    urlcolor=cyan,
}

% ==== Beamer customization section ================

\title{Введение в разработку на C++}
\subtitle{Введение}
\author{Пыхов~Олег~Александрович\inst{1} \\ \texttt{pyhovoleg31@gmail.com}}
\institute{
    \inst{1}%
    ООО~"Прософт-Системы"
}
\date{2024}
\titlegraphic{
    \includegraphics[scale=0.1,valign=t]{Presentations/images/university-logo.png}
    \includegraphics[scale=0.3,valign=t]{Presentations/images/prosoft-logo.png}
}
\logo{
    \begin{tikzpicture}[overlay,remember picture]
        \node[left=0.2cm] at (current page.32){
            \includegraphics[scale=0.15,valign=t]{Presentations/images/prosoft-logo.png}
            \includegraphics[scale=0.05,valign=t]{Presentations/images/university-logo.png}
        };
    \end{tikzpicture}
}

\setbeamertemplate{footline}[frame number]

\begin{document}

\frame[plain]{\titlepage}

\begin{frame}

    \frametitle{Темы курса}

    \begin{itemize}

        \item Компиляция

            \begin{itemize}

                \item Фазы трансляции

                \item Системы сборки

                \item Кросс-компиляция

                \item Управление зависимостями

            \end{itemize}

        \item Система контроля версий (git)

        \item Типы объектов

            \begin{itemize}

                \item Фундаментальные типы

                \item Квалификаторы типов

                \item Операции и операторы

                \item Приведение типов

            \end{itemize}

        \item Условия и циклы

        \item Указатели и работа с памятью

    \end{itemize}

\end{frame}

\begin{frame}

    \frametitle{Темы курса}

    \begin{itemize}

        \item Функции

            \begin{itemize}

                \item Сигнатура функции

                \item Объявление и определение функции

                \item Рекурсия

                \item Перегрузка функций

            \end{itemize}

        \item Пространства имен

        \item Классы

            \begin{itemize}

                \item Поля и методы класса, конструкторы и деструкторы,
                    время жизни объектов, инициализация

                \item \texttt{const}, \texttt{mutable}, статические поля и методы

                \item Наследование

                \item Полиморфизм

            \end{itemize}

    \end{itemize}

\end{frame}

\begin{frame}

    \frametitle{Темы курса}

    \begin{itemize}

        \item Исключения

        \item Шаблоны

        \item Стандартная библиотека

            \begin{itemize}

                \item Строки, потоки ввода-вывода, работа с файлами

                \item Контейнеры

                \item Алгоритмы

                \item Умные указатели

            \end{itemize}

    \end{itemize}

\end{frame}

\begin{frame}[fragile]

    \frametitle{Hello, world!}

    \begin{listing}[H]
        \inputminted[frame=single, fontsize=\footnotesize, linenos]{cpp}
            {Presentations/1-Intro/hello-world.cpp}
        \caption{hello-world.cpp}
    \end{listing}

    \begin{Verbatim}[frame=single, fontsize=\footnotesize, commandchars=\\\{\}]
$ \textbf{g++ hello-world.cpp}
$ \textbf{./a.out}
Hello, World!
    \end{Verbatim}

\end{frame}

\begin{frame}[fragile]

    \frametitle{Стандарты языка C++}

    \begin{columns}

        \begin{column}{0.4\textwidth}

            \begin{itemize}

                \item C++98
                \item C++03
                \item C++11/C++0x
                \item C++14
                \item C++17
                \item C++20
                \item C++23
                \item C++26 (еще не выпущен)

            \end{itemize}

        \end{column}

        \begin{column}{0.6\textwidth}

            \begin{minted}[frame=single, fontsize=\footnotesize]{cpp}
auto i = 3.14; // C++11

// C++14
auto glambda = [] (auto a) {
    return a;
};

// C++17
void f([[maybe_unused]] int a);

// C++20
bool c = a <=> b;

// C++23
std::format("i={}\n", i);
            \end{minted}
        \end{column}

    \end{columns}

\end{frame}

\begin{frame}

    \frametitle{Список полезных источников \\ и ресурсов}

    \begin{enumerate}

        \item \url{https://en.cppreference.com/w/}

        \item \url{https://cplusplus.com/}

        \item Bjarne Stroustrup: The C++ Programming Language

        \item Scott Meyers: Effective Modern C++

        \item \url{https://isocpp.org/}

        \item \href{http://isocpp.github.io/CppCoreGuidelines/CppCoreGuidelines}{C++ Core Guidelines}

        \item Working Draft, Standard for Programming Language C++

        \item \href{https://www.youtube.com/@lectory_fpmi}
            {Лекции преподавателя ФПМИ Ильи Мещерина}

    \end{enumerate}

\end{frame}

\end{document}
